\section{System Design}

The source of truth is Lean plus explicit ledgers.  A mathematical ingredient
is never treated as proved merely because it is plausible in prose.  If it is
not a compiled Lean declaration, it must be a \ProofObligation{} or a cited
result with a recorded Lean status.

The workflow has two public modes.  In \Lean{faithfulPaper} mode, the theorem,
assumptions, constants, and proof target are fixed by the source paper.  In
\Lean{exploratoryProof} mode, the system tracks proof routes for active
research drafts, such as PRIVATE, without promoting candidate assumptions to
facts.

The project-paper framing follows the automation emphasis of
\url{https://github.com/wanshuiyin/Auto-claude-code-research-in-sleep}, the search-memory perspective of \url{https://github.com/Trinkle23897/learning-beyond-gradients}, the
population-and-archive idea of \url{https://github.com/FeiLiu36/EoH}, the Lean-specific blueprint
control loop of \url{https://github.com/YuanheZ/LeanMarathon}, and the Lean proof-diagnostics
workflow exemplified by \url{https://github.com/math-ai-org/mathcode}.  The acceptance condition is
stricter than search success: every cycle must leave a buildable repository or
an explicit blocked status, and post-blueprint cycles must identify the
dynamic leaf or illness area they changed.

\input{figures/agent_stack}
